\documentclass[11pt]{article}
\usepackage[utf8]{inputenc}
\usepackage[margin=1in]{geometry}
\usepackage{enumitem}
\usepackage{hyperref}

\title{Networking Club Activity 2: STP vs. RSTP Performance Analysis}
\author{Networking Club}
\date{March 2, 2026}

\begin{document}

\maketitle

\section{Introduction}
Today we are reviewing a research paper comparing the performance of \textbf{Spanning Tree Protocol (STP)} and \textbf{Rapid Spanning Tree Protocol (RSTP)} via Cisco simulation. This activity aligns with our core club goals: research review, professional certification preparation, and hands-on laboratory experience.

\section{Background: Switching and Loop Prevention}
The primary purpose of a switch is to facilitate data sharing between devices within a Local Area Network (LAN). Switches operate using a specific logic:
\begin{enumerate}
    \item Store the receiver's address.
    \item Check the address table for the receiver.
    \item If the receiver is unknown, \textbf{broadcast} the message to all ports.
\end{enumerate}

In a redundant (loop) topology, these broadcasts can circulate infinitely, resulting in a \textbf{broadcast storm}. To prevent network collapse, we use protocols like STP and RSTP to build a logical "tree" topology, ensuring only one active path exists between any two nodes.

\subsection{The Election Process}
During initialization, switches exchange \textbf{Bridge Protocol Data Units (BPDUs)}. They "shout" their Bridge ID, which is a combination of a configurable \textbf{Priority ID} and the hardware \textbf{MAC ID}. The switch with the lowest value is elected as the \textbf{Root Bridge}—the "brain" of the network.

\subsection{STP vs. RSTP}
\begin{itemize}
    \item \textbf{STP (Legacy):} Uses a one-way communication model. The Root Bridge dictates the topology, and timers are used to detect failures. Convergence is slow because the switch must "wait and think" before transitioning states.
    \item \textbf{RSTP (Modern):} Uses a two-way handshake. Switches actively share topology information. If a "hello" message is missed, neighbors immediately communicate to synchronize a new topology. This results in significantly faster convergence at the cost of constant background communication.
\end{itemize}

\section{Critical Review of Existing Research}
The paper we are reviewing calculates packet loss percentages. However, it suffers from several flaws:
\begin{itemize}
    \item \textbf{Small Sample Size:} Conclusions are based on a limited number of pings.
    \item \textbf{Obvious Findings:} It attempts to "prove" RSTP is better than STP, which is an industry standard (akin to proving a wheel rolls better than a block).
\end{itemize}
\textbf{Note:} This is an example of research that lacks the rigor required for citation. Today, we will fix these flaws by gathering real-world data from physical hardware.

\section{Today's Lab: PVST+ and RPVST+}
In our experiment, we will utilize \textbf{Per-VLAN Spanning Tree (PVST+)} and \textbf{Rapid PVST+}. 
\begin{quote}
    \textit{Note:} Since we are working within a single VLAN today, these will behave like STP and RSTP. A \textbf{VLAN} virtually splits a physical switch into multiple logical switches to contain broadcast domains.
\end{quote}

\section{Collaborative Technical White Paper}
We will synthesize our findings into a benchmarking report. To credit everyone’s work, we are using a tiered contribution model:
\begin{itemize}
    \item \textbf{Technical Contributors:} Execute the lab and submit verified datasets.
    \item \textbf{Lead Authors:} Draft the analysis, results, and narrative.
    \item \textbf{Peer Reviewers:} Edit text and verify calculations via GitHub.
\end{itemize}
The final report will be published on GitHub/Discord, providing a permanent link for your \textbf{Resume or LinkedIn}.

\section{Activity Instructions}
Please refer to the lab guide in the Discord resource section. Advanced challenges include:
\begin{enumerate}
    \item \textbf{Automation:} Automate the initial switch configurations.
    \item \textbf{Optimization:} Analyze the diagram to manually set priorities and compare them against our collected results.
\end{enumerate}

Raquel is available at A and B decks to demonstrate physical wiring and diagram interpretation.

\end{document}